% Unit 085 terjemahan Bahasa Indonesia dari chapter6.tex baris 2458--2603.
% Sumber: Methods of Algebra, Volume 2, commit
% 9a5803ff77dd3257484cb177f851a73770a59dd3 (CC BY 4.0).
% stable-unit-id: o014.aljabr2.chapter6.profinite-cohomology-b
% Cakupan: Bagian 6.12, funktor delta, interpretasi derajat rendah,
% identifikasi sebagai funktor turunan, induksi profinit, dan Lema Shapiro.

% segment-id: o014.li2.chapter6.profinite-cohomology-b.t001
\begin{theorem}\label{prop:profinite-cohomology-long}
	Misalkan $G$ grup profinit dan
	$0\to M'\xrightarrow{f}M\xrightarrow{f'}M''\to0$ barisan eksak pendek
	modul-$G$ mulus. Terdapat barisan eksak panjang terkait
	\begin{multline*}
		\cdots\to\Hm^{n-1}(G,M'')\xrightarrow{\delta^{n-1}}
		\Hm^n(G,M')\xrightarrow{\Hm^n(f)}\Hm^n(G,M)\\
		\xrightarrow{\Hm^n(f')}\Hm^n(G,M'')\xrightarrow{\delta^n}
		\Hm^{n+1}(G,M')\to\cdots,
	\end{multline*}
	dengan konvensi $\Hm^n(G,\cdot)=0$ untuk $n<0$. Homomorfisme penghubung
	$\delta^n$ natural terhadap morfisme barisan eksak pendek; dengan kata lain,
	keluarga ini membentuk funktor-$\delta$ kohomologis.
\end{theorem}
\begin{proof}
	Kita bekerja dengan kompleks standar; argumen untuk versi ternormalisasi
	serupa. Menurut Proposisi~\ref{prop:long-exact-sequence-ses}, cukup buktikan
	bahwa
	\[ 0\to C^n(G,M')\to C^n(G,M)\to C^n(G,M'')\to0 \]
	merupakan barisan eksak pendek. Satu-satunya bagian yang tidak langsung
	adalah kesurjektifan $C^n(G,M)\to C^n(G,M'')$. Misalkan
	$u:G^n\to M''$ kontinu. Seperti di atas, terdapat subgrup normal buka
	$K\lhd G$ sedemikian sehingga $u$ konstan pada setiap koset
	$K^n\lhd G^n$. Pada setiap koset, pilih sembarang prapeta dalam $M$ dari
	nilai konstan tersebut.
\end{proof}

% segment-id: o014.li2.chapter6.profinite-cohomology-b.p001
Ambil modul-$G$ mulus $M$ dengan topologi diskret. Suku berderajat rendah dari
kompleks $C(G,M)$ atau $\overline{C}(G,M)$ masih memiliki interpretasi
langsung.
\begin{description}
	\item[($n=0$)] Jelas bahwa $C^0(G,M)=M$ dan
	$\Hm^0(G,M)=Z^0(G,M)=M^G$. Versi kompleks standar ternormalisasi sama.
	\item[($n=1$)] Unsur $Z^1(G,M)=\overline{Z}^1(G,M)$ adalah homomorfisme
	silang kontinu $G\to M$. Secara ekuivalen, unsur tersebut adalah
	homomorfisme kontinu $\sigma:G\to M\rtimes G$ yang memenuhi
	$\pi\sigma=\identity_G$, seperti dalam
	Proposisi~\ref{prop:crossed-semidirect-product}. Di sini
	$\pi:M\rtimes G\to G$ adalah proyeksi dan $M\rtimes G$ diberi topologi
	hasil kali. Di sisi lain, unsur $B^1(G,M)$ dibangun oleh homomorfisme silang
	berbentuk $f_m:g\mapsto gm-m$; homomorfisme ini otomatis kontinu.

	\item[($n=2$)] Andaikan $M$ hingga. Unsur
	\[ \Hm^2(G,M)\simeq\frac{Z^2(G,M)}{B^2(G,M)}
	\simeq\frac{\overline{Z}^2(G,M)}{\overline{B}^2(G,M)} \]
	berkorespondensi secara bijektif dengan kelas ekuivalensi ekstensi grup profinit
	\[ 1\to M\to E\xrightarrow{\pi}G\to1
	\qquad\text{(barisan eksak pendek grup)}, \]
	dengan aksi konjugasi $G$ pada $M$ berasal dari struktur modul-$G$ pada
	$M$. Dibandingkan dengan versi dalam Definisi~\ref{def:group-ext} dan
	\eqref{eqn:group-ext-diagram}, ``ekstensi grup profinit'' dan masalah
	klasifikasinya juga mensyaratkan:
	\begin{compactitem}
		\item $E$ merupakan grup profinit, $\pi$ kontinu, dan $M$ disertakan
		sebagai subgrup tertutup dari $E$;
		\item ekuivalensi atau isomorfisme ekstensi didefinisikan oleh
		homomorfisme grup kontinu $E\to E'$, dan pembelahan ekstensi juga
		didefinisikan secara topologis.
	\end{compactitem}

	Selain itu, grup automorfisme ekstensi tersebut isomorfik dengan
	$\overline{Z}^1(G,M)$. Buktinya tidak sulit, tetapi memerlukan sedikit
	bahasa topologi dan karena itu diserahkan sebagai latihan bab ini; syarat
	bahwa $M$ hingga bersifat penting.
\end{description}

% segment-id: o014.li2.chapter6.profinite-cohomology-b.p002
Sekarang kita kembali ke funktor turunan kanan
$(\mathrm{R}^n(\cdot)^G)_{n\geq0}$. Kita dapat mengkaji nilainya pada sembarang
kompleks terbatas ke bawah\footnote{Dengan kata lain, mengkaji hiperkohomologi
grup profinit.}, bahkan operasinya dalam kategori turunan. Namun bagian ini
hanya membahas modul-$G$ mulus; sisanya menjadi latihan.

% segment-id: o014.li2.chapter6.profinite-cohomology-b.t002
\begin{proposition}\label{prop:profinite-cohomology-identification}
	Terdapat isomorfisme funktor-$\delta$ kohomologis
	\[ (\Hm^n(G,\cdot))_{n\geq0}
	\simeq(\mathrm{R}^n(\cdot)^G)_{n\geq0}. \]
\end{proposition}
\begin{proof}
	Teorema~\ref{prop:profinite-cohomology-long} menunjukkan bahwa keduanya
	merupakan funktor-$\delta$ kohomologis, dan pada $n=0$ keduanya sama dengan
	funktor invarian $(\cdot)^G$. Menurut
	Proposisi~\ref{prop:erasable-univ-delta}, cukup tunjukkan bahwa
	$\Hm^n(G,\cdot)$ dapat dihapus untuk $n>0$ dalam arti
	Definisi~\ref{def:erasable}. Misalkan $M$ modul-$G$ mulus dan sertakan ke
	dalam modul-$G$ mulus injektif $I$. Untuk setiap subgrup normal buka
	$K\lhd G$, adjungsi \eqref{eqn:invariant-adjunction-profinite} menunjukkan
	bahwa $(\cdot)^K$ mempertahankan objek injektif. Jadi $I^K$ merupakan
	modul-$G/K$ injektif, dan
	\[ n>0\implies\Hm^n(G,I)
	=\varinjlim_K\Hm^n(G/K,I^K)=0. \]
	Sifat dapat dihapus terbukti.
\end{proof}

% segment-id: o014.li2.chapter6.profinite-cohomology-b.p003
Berdasarkan Proposisi~\ref{prop:profinite-cohomology-identification}, metode
\S\ref{sec:change-of-group} dapat diterapkan langsung pada homomorfisme
kontinu $\varphi:H\to G$ dari grup profinit untuk mendefinisikan operasi
``tarikan balik'' kohomologi
\[ \Hm^n(G,M)\to\Hm^n(H,\varphi^*M),\qquad
	M:\text{ modul-$G$ mulus},\quad n\in\Z_{\geq0}. \]
Operasi ini kompatibel dengan barisan eksak panjang. Seperti biasa, gunakan
keeksakan $\varphi^*$ dan sifat funktor-$\delta$ kohomologis universal untuk
mereduksi masalah ke penyertaan yang jelas pada $n=0$,
$M^G\hookrightarrow(\varphi^*M)^H$.

% segment-id: o014.li2.chapter6.profinite-cohomology-b.p004
Sebagai kasus khusus, untuk semua $n$ dan semua subgrup tertutup $H$ terdapat
pemetaan restriksi
$\mathrm{res}^n:\Hm^n(G,M)\to\Hm^n(H,M)$; di sini
$\Res^G_HM$ disingkat sebagai $M$.

% segment-id: o014.li2.chapter6.profinite-cohomology-b.p005
Homomorfisme di atas juga dapat dibangun dari homomorfisme-homomorfisme
\begin{align*}
	\Hm^n(G/K,M^K)&\xrightarrow{\eqref{eqn:change-of-group-coh}}
	\Hm^n(H/\varphi^{-1}(K),\varphi_K^*(M^K))\\
	&\xrightarrow{M^K\subset(\varphi^*M)^{\varphi^{-1}K}}
	\Hm^n(H/\varphi^{-1}(K),(\varphi^*M)^{\varphi^{-1}K}),
\end{align*}
dengan $K$ berkisar atas subgrup normal buka dari $G$. Dengan mereduksi ke
kasus grup hingga, homomorfisme pada kompleks standar dihitung oleh rumus
Proposisi~\ref{prop:grp-equiv}. Untuk membuktikan bahwa pemetaan ini memang
sama dengan pemetaan yang sebelumnya dibangun melalui sifat universal, cukup
bandingkan kasus $n=0$ dan tunjukkan kompatibilitasnya dengan barisan eksak
panjang. Kompleks standar memadai untuk kedua hal tersebut.

% segment-id: o014.li2.chapter6.profinite-cohomology-b.p006
Dengan cara serupa, homomorfisme korestriksi dalam
Definisi--Proposisi~\ref{def:cor} meluas ke grup profinit. Misalkan $H$
subgrup buka dari $G$. Maka $H$ otomatis tertutup dan $(G:H)$ hingga
(lihat \S\ref{sec:profinite-groups}). Definisikan homomorfisme
\[ \mathrm{cor}^n:\Hm^n(H,M)\to\Hm^n(G,M),\qquad
	M:\text{ modul-$G$ mulus},\quad n\in\Z_{\geq0}, \]
yang kompatibel dengan barisan eksak panjang dan pada $n=0$ memberikan
$\nu^{G|H}:M^H\to M^G$.

% segment-id: o014.li2.chapter6.profinite-cohomology-b.p007
Konstruksinya kembali memakai sifat funktor-$\delta$ kohomologis universal,
tetapi di sini kita perlu mengetahui bahwa
$\Res^G_H:G\dcate{Mod}^\infty\to H\dcate{Mod}^\infty$ mempertahankan objek
injektif; ini adalah isi Korolari~\ref{prop:Res-injective-profinite}. Setelah
fakta tersebut diterima, Proposisi~\ref{prop:res-cor} tetap mempunyai analog:
untuk semua $n$,
\[ \mathrm{cor}^n\circ\mathrm{res}^n
	=(G:H)\cdot\identity_{\Hm^n(G,M)}. \]
Seperti sebelumnya, pembuktiannya direduksi ke kasus $n=0$.

% segment-id: o014.li2.chapter6.profinite-cohomology-b.p008
Sekarang kita membahas modul terinduksi dalam kasus profinit. Selanjutnya $G$
selalu merupakan grup profinit. Untuk membedakannya, tuliskan funktor induksi
dalam kasus tanpa topologi (Definisi~\ref{def:induced-module}) sebagai
$\Ind^{G,\mathrm{alg}}_H$; menurut Lema~\ref{prop:Ind-as-mappings}, kita
merealisasikannya sebagai ruang pemetaan.

% segment-id: o014.li2.chapter6.profinite-cohomology-b.d001
\begin{definition}
	\index{modul-$G$!terinduksi (induced)}
	Untuk subgrup tertutup $H$ dari $G$ dan modul-$H$ mulus $N$, lengkapi ruang
	pemetaan $\mathrm{Maps}(G,N):=\{f:G\to N\}$ dengan operasi penjumlahan dan
	perkalian skalar titik demi titik, serta aksi kiri $G$
	\[ (gf)(x)=f(xg),\qquad g,x\in G. \]
	Definisikan modul-$G$ mulus
	\begin{align*}
		\Ind^G_H(N)&:=\{f\in\mathrm{Maps}(G,N)^\infty:
		\forall h\in H,\ \forall x\in G,\ f(hx)=hf(x)\}\\
		&=(\Ind^{G,\mathrm{alg}}_H(N))^\infty.
	\end{align*}
	Notasi $(\cdot)^\infty$ berasal dari
	Definisi--Proposisi~\ref{def:smooth-part}. Objek ini disebut modul-$G$
	terinduksi dari $N$.
\end{definition}

% segment-id: o014.li2.chapter6.profinite-cohomology-b.p009
Karena $G$ kompak, kemulusan $f\in\Ind^{G,\mathrm{alg}}_H(N)$ ekuivalen
dengan kontinuitasnya sebagai pemetaan $G\to N$, dengan $N$ diberi topologi
diskret.

% segment-id: o014.li2.chapter6.profinite-cohomology-b.t003
\begin{proposition}\label{prop:Ind-injective-profinite}
	Untuk setiap subgrup tertutup $H$ dari $G$, terdapat pasangan adjoin
	\[ \Res^G_H:G\dcate{Mod}^\infty
	\leftrightarrows H\dcate{Mod}^\infty:\Ind^G_H. \]
	Sebagai akibatnya, $\Ind^G_H$ mempertahankan objek injektif.

	Jika $H$ subgrup buka dari $G$, terdapat pula pasangan adjoin
	\[ \Ind^G_H:H\dcate{Mod}^\infty
	\leftrightarrows G\dcate{Mod}^\infty:\Res^G_H. \]
	Dalam keadaan ini $\Ind^G_H$ juga mempertahankan objek projektif.
\end{proposition}
\begin{proof}
	Berdasarkan adjungsi dalam kasus tanpa topologi dan
	Definisi--Proposisi~\ref{def:smooth-part}, untuk setiap modul-$G$ mulus $M$
	dan modul-$H$ mulus $N$ terdapat isomorfisme kanonik
	\begin{multline*}
		\Hom_H(\Res^G_H(M),N)
		\simeq\Hom_G(M,\Ind^{G,\mathrm{alg}}_H(N))\\
		=\Hom_G(M,\Ind^{G,\mathrm{alg}}_H(N)^\infty)
		=\Hom_G(M,\Ind^G_H(N)).
	\end{multline*}

	Sekarang andaikan $H$ buka, sehingga $(G:H)$ hingga. Dari
	Korolari~\ref{prop:iInd-Ind} kita memperoleh adjungsi tanpa topologi
	\[ \Hom_H(N,\Res^G_H(M))
	\simeq\Hom_G(\Ind^{G,\mathrm{alg}}_H(N),M). \]
	Kita nyatakan bahwa
	$\Ind^{G,\mathrm{alg}}_H(N)=\Ind^G_H(N)$. Ambil
	$f\in\Ind^{G,\mathrm{alg}}_H(N)$ dan $x\in G$, lalu pilih subgrup buka
	$K_0\subset H$ dengan $f(x)\in N^{K_0}$. Definisikan
	$K:=x^{-1}K_0x$, yang merupakan subgrup buka dari $G$. Maka
	\[ g\in K\implies xg=xgx^{-1}x\in K_0x
	\implies f(xg)=f(x). \]
	Karena $x$ sembarang, pemetaan $f:G\to N$ lokal konstan, sehingga
	$f\in\Ind^G_H(N)$. Pasangan adjoin kedua terbukti.

	Terakhir, pernyataan tentang pemertahanan objek injektif atau projektif
	mengikuti dari keeksakan $\Res^G_H$ dan
	Proposisi~\ref{prop:adjoint-injective-projective}.
\end{proof}

% segment-id: o014.li2.chapter6.profinite-cohomology-b.l001
\begin{lemma}\label{prop:Ind-exact-profinite}
	Jika $H$ subgrup tertutup dari $G$, maka
	$\Ind^G_H:H\dcate{Mod}^\infty\to G\dcate{Mod}^\infty$ eksak.
\end{lemma}
\begin{proof}
	Kita telah mengetahui bahwa $\Ind^G_H$ memiliki adjoin kiri, sehingga
	funktor ini eksak kiri. Tinggal membuktikan bahwa
	$N_1\twoheadrightarrow N_2$ mengakibatkan
	$\Ind^G_H(N_1)\twoheadrightarrow\Ind^G_H(N_2)$. Menurut
	Lema~\ref{prop:profinite-section}, homomorfisme hasil bagi
	$\pi:G\to H\backslash G$ mempunyai penampang kontinu $s$. Untuk setiap
	modul-$H$ mulus $N$, terdapat isomorfisme modul-$\Bbbk$
	\[\begin{tikzcd}[row sep=tiny]
		\Ind^G_H(N) \arrow[leftrightarrow,r,"\sim"] &
		\{\text{pemetaan kontinu }H\backslash G\to N\}\\
		f \arrow[mapsto,r] & {[x\mapsto f(s(x))]}\\
		{[hs(x)\mapsto hf'(x)]} & f' \arrow[mapsto,l]
	\end{tikzcd}\]
	dengan $x\in H\backslash G$ dan $h\in H$. Walaupun isomorfisme ini tidak
	mempertahankan aksi $G$, ia natural dalam $N$.

	Karena pemetaan kontinu $H\backslash G\to N$ selalu lokal konstan, bekerja
	pada ruas kanan segera menunjukkan bahwa
	$\Ind^G_H(N_1)\to\Ind^G_H(N_2)$ surjektif.
\end{proof}

% segment-id: o014.li2.chapter6.profinite-cohomology-b.t004
\begin{corollary}\label{prop:Res-injective-profinite}
	Jika $H$ subgrup buka dari $G$, maka
	$\Res^G_H:G\dcate{Mod}^\infty\to H\dcate{Mod}^\infty$
	mempertahankan objek injektif.
\end{corollary}
\begin{proof}
	Menurut dua hasil sebelumnya, funktor ini memiliki adjoin kiri eksak
	$\Ind^G_H$.
\end{proof}

% segment-id: o014.li2.chapter6.profinite-cohomology-b.p010
Sekarang kita memperoleh versi profinit dari Teorema~\ref{prop:Shapiro}.

% segment-id: o014.li2.chapter6.profinite-cohomology-b.t005
\begin{proposition}
	\index{Lema Shapiro}
	Misalkan $H$ subgrup tertutup dari grup profinit $G$ dan $N$ modul-$H$
	mulus. Terdapat isomorfisme
	\[ \Hm^n(G,\Ind^G_H(N))\simeq\Hm^n(H,N),
	\qquad n\in\Z_{\geq0}. \]
\end{proposition}
\begin{proof}
	Karena $\Ind^G_H$ eksak dan mempertahankan objek injektif
	(Lema~\ref{prop:Ind-exact-profinite} dan
	Proposisi~\ref{prop:Ind-injective-profinite}), argumen dalam
	Teorema~\ref{prop:Shapiro} yang memakai funktor-$\delta$ kohomologis yang
	dapat dihapus berlaku tanpa perubahan. Ingat bahwa inti argumen tersebut adalah
	kasus $n=0$, yakni
	\[ \Ind^G_H(N)^G
	=(\Ind^{G,\mathrm{alg}}_H(N)^\infty)^G
	=(\Ind^{G,\mathrm{alg}}_H(N))^G\rightiso N^H. \]
	Kesamaan pertama adalah definisi, kesamaan kedua mudah, dan isomorfisme
	terakhir berasal dari Korolari~\ref{prop:Ind-invariant}.
\end{proof}

% segment-id: o014.li2.chapter6.profinite-cohomology-b.p011
Selain itu:
\begin{itemize}
	\item barisan spektral kohomologis Lyndon--Hochschild--Serre
	(lihat \S\ref{sec:LHS-SS}) meluas ke keadaan ketika $G$ grup profinit dan
	$H$ subgrup normal tertutup;
	\item hasil kali cawan (lihat \S\ref{sec:cup-product}) juga meluas ke grup
	profinit dengan memakai definisi langsung pada kompleks standar; lihat
	Definisi--Proposisi~\ref{def:group-cup-3}.
	\index{hasil kali cawan}
\end{itemize}
Pernyataannya tidak berubah, sehingga tidak kita ulangi.
